%% =====================================================================
%%  Chapitre 2 — Python as a calculator
%%  Tout le code est vérifié par outils/verifier_code.py (Shell IDLE simulé).
%% =====================================================================
\chapitre{2}{Python as a calculator}{Numbers, operators and expressions}{illusCalcul}

\histoire{Computers were invented to calculate: three centuries ago, a \enquote{computer} was a person
whose job was doing sums by hand. Python inherits this talent. It handles numbers from the tiniest to
the astronomically large, faster than any pocket calculator, and never makes a mistake\,---\,as long as
you tell it exactly what to do. In this chapter, you learn Python's operators, its two kinds of
numbers and the rules that decide in which order a calculation is carried out.}

\begin{objectifs}
\item add, subtract, multiply and divide with Python;
\item tell whole numbers (\py{int}) from decimal numbers (\py{float});
\item use floor division \py{//}, the remainder \py{\%} and powers \py{**};
\item predict the order in which a calculation is done;
\item round, take absolute values, find the smallest and the largest;
\item understand why some decimal results look strange;
\item avoid the classic calculation mistakes.
\end{objectifs}

\auprogramme
  {expression, operator, whole and decimal numbers, order of operations}
  {\py{+ - * /}, \py{// \% **}, \py{round()}, \py{abs()}, \py{min()}, \py{max()}, \py{type()}}
  {the Shell as a pocket calculator}
  {spaces around operators; parentheses for clarity; \py{1\_000\_000}}
  {a holiday budget calculator}

\seances[4]{livre/otNavy/Numbers/The four operations and the two kinds of numbers (2.1--2.2),
  blocs/otViolet/More operators/Floor division and remainder and powers (2.3),
  cible/otTeal/Priorities/Order of operations and useful functions (2.4--2.5),
  crayon/otGreenDark/Practise/Under the hood and mini-project and exercises (2.6--2.11)}

%% ---------------------------------------------------------------------
\rubrique{learn}{The four operations}

Open the IDLE Shell. Python understands the four operations of arithmetic, written with these signs:
\py{+} for addition, \py{-} for subtraction, \py{*} for multiplication and \py{/} for division.

\begin{console}
>>> 7 + 5
12
>>> 7 - 5
2
>>> 7 * 5
35
>>> 7 / 5
1.4
\end{console}

A calculation such as \py{7 + 5} is called an \textbf{expression}: a combination of values and
\textbf{operators} that Python \textbf{evaluates} to produce a single value, its result. The values
on either side of an operator are its \textbf{operands}.

\begin{definition}{expression}
A piece of code that Python evaluates to produce a value: \py{7 + 5}, \py{3 * 4 - 1}, or simply
\py{42}. An \textbf{operator} is a sign that combines values, such as \py{+} or \py{*}.
\end{definition}

\soustitre{Division always gives a decimal number}
Look closely at these two divisions:

\begin{console}
>>> 10 / 2
5.0
>>> 9 / 3
3.0
\end{console}

Even when the division is exact, the operator \py{/} gives a result with a decimal point:
\py{5.0}, not \py{5}. You will see in section~2.2 that \py{5} and \py{5.0} are two different kinds
of numbers for Python.

\soustitre{Negative and very large numbers}
A minus sign in front of a number makes it negative. And Python's whole numbers have no size limit:
you can multiply huge numbers without losing a single digit.

\begin{console}
>>> -5 + 3
-2
>>> 123456789 * 987654321
121932631112635269
\end{console}

\begin{atoi}[Time calculations]
In the Shell, work out the number of minutes in a week ($7$ days of $24$ hours of $60$ minutes), then
the number of seconds in a year of 365 days.
\end{atoi}

%% ---------------------------------------------------------------------
\rubrique{learn}{Two kinds of numbers}

Python distinguishes two kinds of numbers. \textbf{Integers} are whole numbers, positive or negative,
with no decimal point: \py{12}, \py{-3}, \py{2026}. \textbf{Floating-point numbers}, or
\textbf{floats}, have a decimal point: \py{3.5}, \py{-0.25}, \py{5.0}. Each kind of value is a
\textbf{type}, and the function \py{type()} tells you the type of any value:

\begin{console}
>>> type(42)
<class 'int'>
>>> type(3.5)
<class 'float'>
>>> type(10 / 2)
<class 'float'>
\end{console}

\py{int} is short for \emph{integer}, \py{float} for \emph{floating-point number}. Ignore the word
\py{class} for now: it simply means \enquote{type}.

\begin{definition}{type}
The kind of a value, which decides what you can do with it. \py{int} is the type of whole numbers,
\py{float} the type of decimal numbers. The function \py{type()} gives the type of a value.
\end{definition}

\soustitre{Mixing integers and floats}
When a calculation mixes an \py{int} and a \py{float}, the result is a \py{float}, even if it has no
decimal part:

\begin{console}
>>> 2 + 0.5
2.5
>>> 3 * 2.0
6.0
>>> 4 * 25
100
\end{console}

\begin{regle}[int or float?]
\begin{itemize}
\item \py{+}, \py{-} and \py{*} between two integers give an integer;
\item as soon as one operand is a float, the result is a float;
\item the division \py{/} always gives a float.
\end{itemize}
\end{regle}

\soustitre{From one type to the other}
The functions \py{int()} and \py{float()} convert a number to the other type. Be careful:
\py{int()} does not round, it simply \emph{cuts off} the decimal part.

\begin{console}
>>> float(7)
7.0
>>> int(3.9)
3
>>> int(-3.9)
-3
\end{console}

\begin{attention}[A decimal point, never a comma]
In many countries, people write \enquote{3,5} for three and a half. In Python, the decimal separator
is always a \textbf{dot}. With a comma, Python does not report an error: it understands two separate
numbers, and you get a strange result that you will meet again in chapter~19.
\begin{console}
>>> 3,5
(3, 5)
>>> 3,5 * 2
(3, 10)
\end{console}
\end{attention}

%% ---------------------------------------------------------------------
\rubrique{learn}{Floor division, remainder and powers}

Imagine you share 17 sweets between 5 children. Each child gets 3 sweets, and 2 sweets are left over.
Python has an operator for each of these two numbers: \py{//} gives the \textbf{quotient} of the
\textbf{floor division} (3), and \py{\%} gives the \textbf{remainder} (2).

\begin{center}
\begin{tikzpicture}[x=1mm,y=1mm]
  \foreach \k in {0,...,4}{%
    \fill[otNavyTint,rounded corners=1.6mm] (\k*24,0) rectangle ++(20,13);
    \foreach \j in {0,1,2}{\fill[otNavy] (\k*24+4.5+\j*5.5,6.5) circle (1.9);}
    \node[font=\otMedium\fontsize{7.4}{9}\selectfont,text=otMuted] at (\k*24+10,-3.2) {child \the\numexpr\k+1\relax};}
  \foreach \j in {0,1}{\fill[otOrange] (126+\j*5.5,6.5) circle (1.9);}
  \draw[otOrange,line width=0.6pt,rounded corners=1.6mm,dashed] (122,0) rectangle ++(14,13);
  \node[font=\otMedium\fontsize{7.4}{9}\selectfont,text=otOrangeDark] at (129,-3.2) {left over};
  \node[anchor=west,font=\ttfamily\fontsize{9.5}{11}\selectfont,text=otNavy] at (141,9) {17 // 5 = 3};
  \node[anchor=west,font=\ttfamily\fontsize{9.5}{11}\selectfont,text=otOrangeDark] at (141,3) {17 \% 5 = 2};
\end{tikzpicture}
\end{center}

\begin{console}
>>> 17 // 5
3
>>> 17 % 5
2
>>> 5 * 3 + 2
17
\end{console}

The last line checks the result: 5 children times 3 sweets, plus the 2 left over, gives back the 17
sweets. This check always works: the divisor times the quotient, plus the remainder, gives back
the number you started from.

\soustitre{What are they good for?}
These two operators are surprisingly useful. To convert 135 minutes into hours and minutes, divide
by 60: the quotient gives the hours, the remainder the minutes.

\begin{console}
>>> 135 // 60
2
>>> 135 % 60
15
>>> print(135 // 60, "h", 135 % 60, "min")
2 h 15 min
\end{console}

The remainder of a division by 2 tells you whether a number is even (remainder 0) or odd
(remainder~1). And \py{//} and \py{\%} by 10 separate the last digit of a number from the others:

\begin{console}
>>> 37 % 2
1
>>> 1234 % 10
4
>>> 1234 // 10
123
\end{console}

\begin{attention}[Floor division rounds down]
\py{//} rounds the result \emph{down}, towards smaller numbers. For positive numbers, it simply drops
the decimal part. For negative numbers, it goes one step further: \py{-7 // 2} gives \py{-4}, not
\py{-3}, because $-4$ is smaller than $-3.5$. You will rarely need it with negative numbers.
\end{attention}

\soustitre{Powers}
The operator \py{**} raises a number to a power: \py{2 ** 3} means $2^3 = 2 \times 2 \times 2$. A
power of \py{0.5} gives the square root.

\begin{console}
>>> 2 ** 3
8
>>> 10 ** 6
1000000
>>> 2 ** 0.5
1.4142135623730951
>>> 2 ** 100
1267650600228229401496703205376
\end{console}

\begin{definition}{remainder}
What is left over after a floor division. In Python, \py{a \% b} is the remainder of the division of
\py{a} by \py{b}, and \py{a // b} its quotient. The sign \py{\%} is also called \enquote{modulo}.
\end{definition}

\begin{atoi}[Minutes and seconds]
A video lasts 1000 seconds. Use \py{//} and \py{\%} to display its length in minutes and seconds,
in the form \py{16 min 40 s}.
\end{atoi}

%% ---------------------------------------------------------------------
\rubrique{learn}{Order of operations}

What is \py{2 + 3 * 4}? If you calculate from left to right, you find 20. Python finds 14:

\begin{console}
>>> 2 + 3 * 4
14
>>> (2 + 3) * 4
20
\end{console}

Like in mathematics, multiplication has \textbf{priority} over addition: Python first works out
\py{3 * 4}, then adds \py{2}. To impose another order, use parentheses: what is between parentheses
is always calculated first.

\begin{tcolorbox}[otBase,colback=white,colframe=otRule,boxrule=0.5pt,left=4mm,right=4mm]
\otEtiquette{otNavy}{Priority of the operators, from highest to lowest}\par\vspace{1.6mm}
\renewcommand{\arraystretch}{1.32}%
\begin{tabularx}{\linewidth}{@{}>{\centering\arraybackslash}p{8mm}p{28mm}Y@{}}
\textbf{1} & \py{( )} & parentheses: always first\\
\textbf{2} & \py{**} & powers\\
\textbf{3} & \py{-x} & the minus sign of a negative number\\
\textbf{4} & \py{*\ \ /\ \ //\ \ \%} & multiplication, divisions, remainder\\
\textbf{5} & \py{+\ \ -} & addition, subtraction\\
\end{tabularx}\par\vspace{1.4mm}
{\fontsize{9.8}{12.4}\selectfont Operators with the same priority are evaluated \textbf{from left to
right}, except powers, which go from right to left.}
\end{tcolorbox}

Here is how Python evaluates a longer expression, step by step. At each step, it performs the
operation with the highest priority (highlighted):

\begin{calcul}
\pas[]{\py{2 + 3 * \hlc{4 ** 2}}}{the power comes first}
\pas{\py{2 + \hlc{3 * 16}}}{then the multiplication}
\pas{\py{\hlc{2 + 48}}}{then the addition}
\pas{\py{50}}{the result}
\end{calcul}

\begin{console}
>>> 2 + 3 * 4 ** 2
50
>>> 20 / 4 * 2
10.0
>>> 2 ** 3 ** 2
512
\end{console}

\py{20 / 4 * 2} is evaluated from left to right: \py{20 / 4} gives \py{5.0}, then \py{5.0 * 2}
gives \py{10.0}. But \py{2 ** 3 ** 2} is evaluated from right to left: \py{3 ** 2} gives \py{9},
then \py{2 ** 9} gives \py{512}.

\begin{attention}[The minus sign and powers]
\py{-2 ** 2} gives \py{-4}, not \py{4}: the power has priority over the minus sign, so Python
calculates $-(2^2)$. Write \py{(-2) ** 2} to square $-2$.
\end{attention}

\begin{atoi}[The average]
Ana's marks are 12, 15 and 18. Calculate their average in the Shell. Then try the same calculation
without parentheses: what do you get, and why?
\end{atoi}

%% ---------------------------------------------------------------------
\rubrique{learn}{Useful functions for numbers}

Python provides several ready-made functions for numbers. You call them like \py{print()}: their
name, then their arguments between parentheses, separated by commas.

\begin{console}
>>> round(3.14159, 2)
3.14
>>> round(2.7)
3
>>> abs(-7)
7
>>> min(4, 9, 2)
2
>>> max(4, 9, 2)
9
\end{console}

\begin{itemize}
\item \py{round(x, n)} rounds \py{x} to \py{n} digits after the decimal point; without \py{n}, it
  rounds to the nearest integer and gives an \py{int};
\item \py{abs(x)} gives the \textbf{absolute value} of \py{x}: the number without its sign;
\item \py{min()} and \py{max()} give the smallest and the largest of their arguments, however many
  there are.
\end{itemize}

You can use the result of a function inside an expression, or even inside another function call.
Python then evaluates from the inside out:

\begin{console}
>>> round(10 / 3, 2)
3.33
>>> max(3, 8) - min(3, 8)
5
>>> print("Rounded:", round(2 ** 0.5, 3))
Rounded: 1.414
\end{console}

\begin{attention}[Rounding exactly halfway]
When a number is exactly halfway between two integers, \py{round()} chooses the \textbf{even}
neighbour: \py{round(1.5)} and \py{round(2.5)} both give \py{2}. This is called \enquote{banker's
rounding}: on average, it does not push results up or down. It rarely matters, but it surprises
everybody the first time.
\begin{console}
>>> round(1.5)
2
>>> round(2.5)
2
\end{console}
\end{attention}

%% ---------------------------------------------------------------------
\rubrique{hood}{Why \texorpdfstring{\py{0.1 + 0.2}}{0.1 + 0.2} is not exactly 0.3}

Type this in the Shell, and get ready for a surprise:

\begin{console}
>>> 0.1 + 0.2
0.30000000000000004
>>> 3 * 2.4
7.199999999999999
\end{console}

Python is not broken. The problem comes from the way computers store decimal numbers. Inside the
machine, everything is written in \textbf{binary}, with only the digits 0 and 1. In binary, most
decimal numbers have no exact representation, just as $\frac{1}{3}$ has no exact representation in
decimal: \py{0.333333}\ldots{} goes on for ever.

\begin{center}
\begin{tikzpicture}[x=1mm,y=1mm]
  \fill[otVioletTint,rounded corners=2mm] (0,0) rectangle (79,22);
  \node[anchor=north west,font=\otSemi\fontsize{8.4}{10}\selectfont,text=otViolet] at (4,19) {In decimal (base 10)};
  \node[anchor=west,font=\ttfamily\fontsize{9.4}{11}\selectfont] at (4,9) {1 / 3 = 0.3333333333\ldots};
  \fill[otVioletTint,rounded corners=2mm] (89,0) rectangle (170,22);
  \node[anchor=north west,font=\otSemi\fontsize{8.4}{10}\selectfont,text=otViolet] at (93,19) {In binary (base 2)};
  \node[anchor=west,font=\ttfamily\fontsize{9.4}{11}\selectfont] at (93,9) {0.1 = 0.000110011001100\ldots};
\end{tikzpicture}
\end{center}

A float keeps about 16 significant digits. Beyond that, the number is cut off, so \py{0.1} is stored
as a value extremely close to 0.1, but not exactly equal. Most of the time, Python hides this tiny
difference when it displays the number; sometimes, after a calculation, the difference shows in the
last digit.

\begin{regle}[Living with floats]
\begin{itemize}
\item Integers are \textbf{always exact}, however large they are.
\item Floats are \textbf{very close approximations}: perfect for measurements, prices, physics.
\item To display a float neatly, round it: \py{round(3 * 2.4, 2)} gives \py{7.2}.
\end{itemize}
\end{regle}

\begin{culture}[A worldwide standard]
Almost all computers in the world store floats in the same way, defined by a standard called
\emph{IEEE~754}. You would get exactly the same \py{0.30000000000000004} in JavaScript, Java or C:
it is the price of fast calculations on very large and very small numbers.
\end{culture}

%% ---------------------------------------------------------------------
\rubrique{debug}{Calculation mistakes}

Some calculation mistakes stop your program with an error message. Others are more treacherous: the
program runs, but the result is wrong.

\cas{1}{Dividing by zero}
Dividing by zero is impossible, in Python as in mathematics. Python stops with a
\py{ZeroDivisionError}, for \py{/} but also for \py{//} and \py{\%}:
\begin{console}
>>> 10 / 0
Traceback (most recent call last):
  File "<pyshell#0>", line 1, in <module>
    10 / 0
ZeroDivisionError: division by zero
\end{console}

\cas{2}{The wrong multiplication sign}
The sign of multiplication is the star \py{*}. Neither the letter \py{x} nor the sign \py{×} work:
\begin{console}
>>> 5 x 3
SyntaxError: invalid syntax
>>> 5 × 3
SyntaxError: invalid character '×' (U+00D7)
\end{console}

\cas{3}{The hat for powers}
In many programs, such as spreadsheets, \py{2\textasciicircum{}3} means \enquote{2 to the power of 3}. Not in
Python! The hat \py{\textasciicircum{}} is an operator on the binary digits of numbers, which you will never need
in this book. No error message, but a wrong result:
\begin{console}
>>> 2 ^ 3
1
>>> 2 ** 3
8
\end{console}

\cas{4}{Forgotten parentheses}
The most frequent mistake of all, because Python cannot detect it: the calculation is valid, it just
does not mean what you wanted.
\begin{console}
>>> 12 + 15 + 18 / 3
33.0
>>> (12 + 15 + 18) / 3
15.0
\end{console}

\begin{pratique}[Check your results]
For every calculation, estimate the order of magnitude before running it. An average of three marks
out of 20 that gives 33 should immediately make you suspicious. A program that runs without an error
message is not necessarily a correct program.
\end{pratique}

%% ---------------------------------------------------------------------
\rubrique{pro}{Readable calculations}

\soustitre{Spaces around operators}
\py{7+5*2} and \py{7 + 5 * 2} give the same result, but the second is much easier to read. The style
guide PEP~8 recommends one space on each side of the operators \py{+}, \py{-}, \py{*}, \py{/},
\py{//}, \py{\%}, and no space just inside parentheses: \py{(2 + 3)}, not \py{( 2 + 3 )}.

\soustitre{Parentheses for clarity}
Even when they are not required, parentheses can make your intention clearer. Python ignores them
when they change nothing, but the reader appreciates them:

\begin{code}{file=clarity.py}
print("Total:", (3 * 4.5) + (2 * 1.75))
print("Total:", 3 * 4.5 + 2 * 1.75)
\end{code}
\begin{sortie}
Total: 17.0
Total: 17.0
\end{sortie}

\soustitre{Underscores in large numbers}
In a number, Python allows underscores \py{\_} between digits, to group them by three like the
spaces of everyday writing. They change nothing to the value:

\begin{console}
>>> 1_000_000 * 3
3000000
\end{console}

\begin{pratique}[Say what the numbers mean]
A bare number in a program, such as \py{0.2} or \py{86400}, is called a \textbf{magic number}:
nobody knows where it comes from. Add a short comment with its meaning and its unit:
\py{86400 \# seconds in a day}. In chapter~3, variables will give these numbers a name.
\end{pratique}

%% ---------------------------------------------------------------------
\rubrique{build}{Mini-project: a holiday budget}

Ana and her friend Leo are planning a week in Lisbon. Here are the costs they found:

\begin{tcolorbox}[otBase,colback=otGreenTint,colframe=otGreenTint,left=5mm]
\begin{itemize}
\item return flight: 89.50 euros per person;
\item hotel: 74 euros per night for the room, 6 nights;
\item food: 25 euros per person per day, for 7 days;
\item museum pass: 21.50 euros per person.
\end{itemize}
\end{tcolorbox}

They want a program that displays each item of the budget, the total, the cost per person, and how
much they must save each month during the 6 months before leaving.

\soustitre{Step 1: one line per item}
Each item is a \py{print()} with a label and a calculation. For the flights, there are 2 people at
89.50 euros each: \py{2 * 89.5}.

\soustitre{Step 2: the total and the shares}
The total adds up the four calculations. Then it is divided by 2 for the cost per person, and by 6
for the monthly savings, rounded to the cent with \py{round()}.

\begin{code}{file=holiday_budget.py}
# holiday_budget.py - a week in Lisbon for 2 people
print("=== HOLIDAY IN LISBON ===")
print("Flights:", 2 * 89.5)
print("Hotel:", 6 * 74)
print("Food:", 2 * 7 * 25)
print("Museums:", 2 * 21.5)
print("Total:", 2 * 89.5 + 6 * 74 + 2 * 7 * 25 + 2 * 21.5)
print("Per person:", (2 * 89.5 + 6 * 74 + 2 * 7 * 25 + 2 * 21.5) / 2)
print("To save each month:",
      round((2 * 89.5 + 6 * 74 + 2 * 7 * 25 + 2 * 21.5) / 6, 2))
\end{code}
\begin{sortie}
=== HOLIDAY IN LISBON ===
Flights: 179.0
Hotel: 444
Food: 350
Museums: 43.0
Total: 1016.0
Per person: 508.0
To save each month: 169.33
\end{sortie}

The last instruction is too long for one line, so it continues on the next one. Python allows this
as long as a parenthesis is still open: it knows that the instruction is not finished.

\soustitre{Step 3: what if Leo's sister comes too?}
With 3 people, the flights, the food and the museums change, and so does every total. Count how many
numbers you would have to change in the program: the \py{2} appears seven times! Forget a single one,
and the budget is wrong. This is exactly the problem that \textbf{variables} solve, in the next
chapter.

\begin{atoi}[Your own trip]
Adapt the program to a trip of your choice: destination, number of people, prices and number of
nights. Check every result with an estimate.
\end{atoi}

%% ---------------------------------------------------------------------
\rubrique{practice}{Exercises}

\exercice[type=predict,niveau=1]{What is the result?}{Without using Python, write the result of each
expression, paying attention to the type: write \py{6.0} if the result is a float. Then check in the
Shell.}
\begin{questions}[4]
\item \py{7 + 3 * 2}
\item \py{(7 + 3) * 2}
\item \py{9 / 3}
\item \py{9 // 2}
\item \py{9 \% 2}
\item \py{2 ** 4}
\item \py{10 - 2 - 3}
\item \py{2 * 3.0}
\end{questions}

\exercice[type=predict,niveau=1]{int or float?}{For each expression, say whether Python gives an
\py{int} or a \py{float}, then check with \py{type()}.}
\begin{questions}[4]
\item \py{5}
\item \py{5.0}
\item \py{3 + 4}
\item \py{3 + 4.0}
\item \py{10 / 5}
\item \py{10 // 5}
\item \py{2 ** 3}
\item \py{round(2.6)}
\end{questions}

\exercice[type=shell,niveau=1]{Shopping in the Shell}{Use the Shell to answer each question. Write
the expression you typed and the result.}
\begin{questions}
\item Five notebooks at 3 euros and two pens at 1.20 euros: how much in total?
\item A 2-litre bottle is shared equally between 6 glasses. How much in each glass, rounded to the
  hundredth of a litre?
\item A jacket at 80 euros has a 30\,\% discount. What is its new price?
\item What is the difference between the most expensive and the cheapest of these prices: 12.5,
  8.75, 19.9, 4.6? Use \py{max()} and \py{min()}.
\end{questions}

\exercice[type=write,niveau=1]{Hours and minutes}{A film lasts 200 minutes. Write a program that
displays its length in hours and minutes, exactly as below. Python must do the calculations.}
\cible{film}
\begin{sortie}
The film lasts 3 h 20 min.
\end{sortie}

\exercice[type=complete,niveau=2]{The missing operators}{Replace each blank with one operator
(\py{+}, \py{-}, \py{*}, \py{/}, \py{//}, \py{\%} or \py{**}) so that the program displays the output
below.}
\begin{code}{file=operators.py, trous}
print(7 ___ 2)
print(7 ___ 2)
print(7 ___ 2)
print(7 ___ 2)
print(7 ___ 2)
\end{code}
\cible{operateurs}
\begin{sortie}
3.5
3
1
49
14
\end{sortie}

\exercice[type=fix,niveau=2]{Wrong but no error}{This program runs without any error message, but
its three results are wrong. Find the three mistakes and correct them so that the program displays
the expected output.}
\begin{code}{file=wrong_results.py, bug}
print("Square of 5:", 5 ^ 2)
print("Half of 7.5:", 7,5 / 2)
print("Average of 12 and 14:", 12 + 14 / 2)
\end{code}
\begin{sortie}
Square of 5: 7
Half of 7.5: 7 2.5
Average of 12 and 14: 19.0
\end{sortie}
\cible{faux}
\begin{sortie}
Square of 5: 25
Half of 7.5: 3.75
Average of 12 and 14: 13.0
\end{sortie}

\exercice[type=write,niveau=2]{Degrees}{To convert a temperature from degrees Celsius to degrees
Fahrenheit, multiply by 9, divide by 5, then add 32. To go back, subtract 32, multiply by 5 and
divide by 9. Write a program that displays the two conversions below.}
\cible{degres}
\begin{sortie}
25 degrees Celsius = 77.0 degrees Fahrenheit
98.6 degrees Fahrenheit = 37.0 degrees Celsius
\end{sortie}

\exercice[type=read,niveau=2]{Step by step}{For each expression, write the steps of the calculation
in the order Python carries them out, then the result.}
\begin{questions}[2]
\item \py{2 + 3 * 4 - 1}
\item \py{(2 + 3) * (4 - 1)}
\item \py{2 ** 2 * 3}
\item \py{20 / 2 / 5}
\item \py{10 - 4 // 3}
\item \py{-3 ** 2 + 10}
\end{questions}

\exercice[type=write,niveau=2]{Splitting the bill}{Seven friends have dinner. The bill is 94.50 euros
and they add a 10\,\% tip. Write a program that displays the total with the tip and the amount each
friend pays, rounded to the cent.}
\cible{addition}
\begin{sortie}
Total with tip: 103.95 euros
Each friend pays: 14.85 euros
\end{sortie}

\exercice[type=explain,niveau=2]{Digit by digit}{Predict the result of each expression, then explain
in one sentence what the operators \py{//} and \py{\%} do with the digits of 1234.}
\begin{questions}[4]
\item \py{1234 \% 10}
\item \py{1234 // 10}
\item \py{1234 // 10 \% 10}
\item \py{1234 // 1000}
\end{questions}

\exercice[type=challenge,niveau=3]{100,000 seconds}{How long are 100,000 seconds in days, hours,
minutes and seconds? Write a program that displays the answer as below. Python must do every
calculation from the number 100000, using \py{//} and \py{\%}.}
\cible{secondes}
\begin{sortie}
100000 seconds = 1 d 3 h 46 min 40 s
\end{sortie}

\exercice[type=challenge,niveau=3]{The legend of the chessboard}{According to an old legend, the
inventor of chess asked to be paid in rice: one grain on the first square of the chessboard, two on
the second, four on the third, and so on, doubling each time up to the 64th square. In total, this
makes $2^{64} - 1$ grains. Write a program that displays the number of grains on the last square,
the total number of grains, and the mass of this rice in tonnes, rounded to the tonne, if a grain
weighs 0.025~g (one tonne is 1,000,000~g).}
\cible{riz}
\begin{sortie}
Grains on the last square: 9223372036854775808
Total number of grains: 18446744073709551615
Total mass in tonnes: 461168601843
\end{sortie}

%% ---------------------------------------------------------------------
\rubrique{review}{Chapter review}

\begin{essentiel}
\otEssTitre{Operators}
\begin{tabularx}{\linewidth}{@{}p{28mm}Y@{}}
\py{+  -  *  /} & the four operations; \py{/} always gives a float\\
\py{//} and \py{\%} & quotient and remainder of the floor division: \py{17 // 5} is 3,
  \py{17 \% 5} is 2\\
\py{**} & power: \py{2 ** 3} is 8; \py{x ** 0.5} is the square root\\
\end{tabularx}

\otEssTitre{Two kinds of numbers}
\py{int}: whole numbers, exact and unlimited. \py{float}: decimal numbers (with a \textbf{dot}),
very close approximations. An operation that involves a float gives a float. \py{type()} gives the
type of a value; \py{int()} cuts off the decimal part.

\otEssTitre{Order of operations}
Parentheses, then \py{**}, then \py{*}~\py{/}~\py{//}~\py{\%}, then \py{+}~\py{-}; from left to
right for equal priorities. When in doubt, add parentheses.

\otEssTitre{Useful functions}
\begin{code}{mini}
round(3.14159, 2)   # 3.14
abs(-7)             # 7
min(4, 9, 2)        # 2
max(4, 9, 2)        # 9
\end{code}

\otEssTitre{Mistakes to avoid}
Dividing by zero (\py{ZeroDivisionError}); \py{x} or \py{\textasciicircum{}} instead of \py{*} or \py{**}; a comma
instead of a decimal point; forgotten parentheses. Estimate every result before trusting it.
\end{essentiel}

\begin{quiz}
\qcm{What does \py{7 / 2} give?}{\py{3}}{\py{3.5}}{\py{3.0}}
\qcm{What does \py{7 // 2} give?}{\py{3}}{\py{3.5}}{\py{1}}
\qcm{What does \py{7 \% 2} give?}{\py{3}}{\py{0.5}}{\py{1}}
\qcm{What is the type of \py{4 * 2.5}?}{\py{int}}{\py{float}}{it depends on the result}
\qcm{What does \py{2 + 3 * 2} give?}{\py{10}}{\py{8}}{\py{12}}
\qcm{How do you write $5^2$ in Python?}{\py{5 ** 2}}{\py{5 \textasciicircum{} 2}}{\py{5 x 2}}
\qcm{What does \py{round(4.567, 1)} give?}{\py{4.5}}{\py{4.6}}{\py{5}}
\qcm{What does \py{int(9.99)} give?}{\py{10}}{\py{9}}{\py{9.0}}
\qcm{Which number is the remainder of 23 divided by 4?}{\py{23 // 4}}{\py{23 \% 4}}{\py{23 / 4}}
\qcm{What does \py{10 / 0} produce?}{\py{0}}{\py{inf}}{a \py{ZeroDivisionError}}
\end{quiz}

\begin{jesais}
\item calculate with \py{+}, \py{-}, \py{*}, \py{/} and \py{**};
\item use \py{//} and \py{\%} to find a quotient and a remainder;
\item tell an \py{int} from a \py{float} and predict the type of a result;
\item predict the order of a calculation, and use parentheses to change it;
\item use \py{round()}, \py{abs()}, \py{min()} and \py{max()};
\item explain why \py{0.1 + 0.2} does not display \py{0.3};
\item recognise a \py{ZeroDivisionError} and the classic calculation mistakes.
\end{jesais}
